% Part: many-valued-logic
% Chapter: supervaluationist-logic
% Section: semantics

\documentclass[../../../include/open-logic-section]{subfiles}

\begin{document}

\olfileid{mvl}{sup}{det}

\olsection{Determinacy operator}

We can introduce operators to talk about whether a sentence is $\True$, 
$\Undef$ or $\False$ in supervaluationism. We extend our language 
with unary !!{operator}s $\Delta$ and $\nabla$.

\begin{defn}[Determinacy operators language]
We add the following clause to our definition of the language:
\begin{itemize}
\item  For every !!{formula} $!A$, $\Delta !A$ and $\nabla !A$ are 
	!!{formula}s too.
\end{itemize}
\end{defn}

\begin{explain}
$\Delta !A$ is meant to capture the idea that $!A$ is \emph{true} (or 
\emph{determinately true} \emph{supertrue}). In supervaluationist 
semantics ,  this means being true at every supervaluation. $\nabla
!A$ is meant to capture the idea that $!A$ is \emph{neither true nor
false} (or  \emph{indeterminate}). In supervaluationist semantics, 
this means being true at some supervaluations but false at others.

Recall that the supervaluationist semantics says that !!a{formula}
$!A$ is $\True$ in !!a{valuation} \pAssign{v} (or !!a{structure} 
\Struct{M}) iff $!A$ is $\True$ in any supervaluation $\pAssign{v^+}$
for $\pAssign{v}$ (or any supervaluation $\Struct{M^+}$ for 
$\Struct{M}$). We want to keep that even if $!A$ is of the form 
$\Delta !B$ or $\nabla !B$. (This will allow us to handle arbitrary 
embeddings of these operators, e.g. $\Delta (\Delta \Obj{p} \lif p)$.)
So we need to assign values of $\Delta !B$ and $\nabla !B$ !!{formula}s
\emph{at a supervaluation}. This is what we do in the semantics below.

It is easy to see that if 
$\Delta !B$ (resp. $\nabla !B$) is true at a supervaluation for some
!!{valuation} \pAssign{v}, it will be true at any other supervaluation
for that !!{valuation}, so the relativisation might seem redundant; but
it is needed keep our general definition of truth as truth at all
supervaluations. Moreover, there are more complex versions of 
supervaluationism that allow $\Delta !B$ !!{formula}s themselves
to vary in truth from one supervaluation to the other. So it is 
good to define their truth relative to supervaluations in this simple
case too, even though the relativization turns out to be redundant in
our simple semantics.
\end{explain}

\begin{defn}[Semantics for determinacy operators]

For propositional logic, let $\pAssign{v^+}$ be a 
supervaluation for $\pAssign{v}$ and $!A$ any !!{formula}:

\begin{itemize}
	\item $\pValue{v^+}(!A)$ is given as in $\LogCL$ if $!A$ is a !!{propositional
	variable} or of the forms $\lnot !B$, $!B \land !C$, $!B \lor !C$, $!B \lif !C$.
	\item \indcase{!A}{\Delta !B}{$\pValue{v^+}(\Delta !B)=\True$ if $\pValue{v_i^+}(!B)=\True$ for any supervaluation $\pAssign{v_i^+}$ for $\pAssign{v}$.
	$\pValue{v^+}(\Delta !B)=\False$ otherwise.
	\item \indcase{!A}{\nabla !B}{$\pValue{v^+}(\Delta !B)=\True$ if $\pValue{v_i^+}(!B)=\True$ and $\pValue{v_j^+}(!B)=\True$ for some supervaluations $\pAssign{v_i^+}$, $\pAssign{v_j^+}$ for $\pAssign{v}$. $\pValue{v^+}(\nabla !B)=\False$ otherwise.
\end{itemize}

For predicate logic, the definition is analogue. Let $\Struct{M^+}$ be a supervaluation of $\Struct{M}$, $s$ an assignment in $\Struct{M^+}$
 and $!A$ any !!{formula}:

\begin{itemize}
	\item $\Value{M^+}(!A)[s]$ is given as in classical logic if $!A$ is of
	of the form $\Atom{f}{t_1, \dots, t_k}$, $\eq[t_1][t_2]$, $\lnot !B$, $!B \land !C$, $!B \lor !C$, $!B \lif !C$, $\lforall[x][!A]$, $\lexists[x][!A]$.
	\item \indcase{!A}{\Delta !B}{$\Value{\Delta !B}{M^+}[s]=\True$ if 
	$\Value{!B}{M_i^+}[s]=\True$ for any supervaluation $\Struct{M_i^+}$ for $\Struct{M}$. $\Value{\Delta !B}{M^+}[s]=\False$ otherwise.
	\item \indcase{!A}{\nabla !B}{$\Value{\Delta !B}{M^+}[s]=\True$ if 
	$\Value{!B}{M_i^+}[s]=\True$ and $\Value{!B}{M_j^+}[s]=\False$ for some
	 supervaluations $\Struct{M_i^+}$, $\Struct{M_j^+}$ for $\Struct{M}$.
	 $\Value{\nabla !B}{M^+}[s]=\False$ otherwise.
\end{itemize}

\end{defn}

Note that the clauses ensure that at any supervaluation $\pAssign{v^+}$,
$\pValue{v^+}(\Delta !A)$ and $\pValue{v^+}(\nabla !A)$ are $\True$
or $\False$, so supervaluations remain bivalent. 

We can use the same clauses for adding determinacy operators 
in classical logic, treating each classical !!{valuation} as its unique own supervaluation.

\begin{prop}[Interdefinability of $\Delta$ and $\nabla$]
$\Delta$ and $\nabla$ are interdefinable. 
\begin{itemize}
	\item $\nabla !A$ is equivalent to $\lnot \Delta !A 
	\land \lnot \Delta \lnot !A$.
	\item $\Delta !A$ is equivalent to $p \land \lnot \nabla p$.
\end{itemize}
\end{prop}

Some consequences. 

\begin{enumerate}
\item $\Entails[\LogSup] \Delta !A \lif !A$
\item $\Delta !A \Entails[\LogSup] !A$
\item $!A \Entails[\LogSup] \Delta !A$. For suppose $\pValue{v}[!A]=
	\True$ for some $\LogSup$ !!{valuation} $\pAssign{v}$. Then 
	$\pValue{v^+}[!A]=\True$ for any supervaluation $\pAssign{v^+}$ 
	for $\pAssign{v}$. So $\pValue{v}[\Delta !A]=
	\True$.
\item $\Entails/[\LogSup] !A \lif \Delta !A$. Counterexample:
	consider $\pAssign{v}(\Obj p)=\Undef$. We have $\pValue{v^+}(\Delta 
	\Obj{p})=\False$
	for any supervaluation $\pAssign{v^+}$ for $\pAssign{v}$, but 
	$\pValue{v_i^+}(\Obj p)=\True$ and $\pValue{v_j^+}(\Obj p)=\False$
	for some supervaluations $\pAssign{v_i^+}$ and $\pAssign{v_i^+}$
	for $\pAssign{v}$, so $\pValue{v}(\Obj{p \lif \Delta p})=\False$.
\item $\lnot \Delta !A \Entails/[\LogSup] \lnot !A$. Counterexample:
	consider $\pAssign{v}(\Obj p)=\Undef$, we have $\pValue{v}(\lnot
	\Delta \Obj{p})=\True$ but $\pValue{v}(\Obj{\lnot p})=\Undef$. 
\item $\Entails[\LogSup] !A \lor \lnot !A$. $\Entails/[\LogSup] \Delta
	!A \lor \Delta \lnot !A$.
\item The rule of \emph{conditional proof} fails in $\LogSup$ 
	with determinacy operators. Conditional proof
	is the rule: if $\Gamma,!A\Entails !B$ then $\Gamma\Entails !A 
	\lif !B$. Consequences 3 and 4 above show that it fails here.
\item The rule of \emph{contraposition} fails in $\LogSup$ 
	with determinacy operators. The contraposition rule is
	the rule: if $\Gamma,!A\Entails !B$ then $\Gamma, \lnot !B\Entails 
	\lnot !A$. Consequences 3 and 5 above show that it fails here. 
\item The rule of \emph{reasoning by cases} fails in $\LogSup$ 
	with determinacy operators. The reasoning by cases rule is 
	the rule: if  $\Gamma\Entails !A \lor !B$ and $\Gamma, !A \Entails !C$
	and $\Gamma, !B \Entails !C$, then $\Gamma\Entails !C$. 
\end{enumerate}
	
\end{document}
